\section{Anthropic、利用規約を改定}
Anthropicは利用規約を1年以上ぶりに更新し、2026年11月12日から適用すると説明した。モデルに対する「持続的かつ不必要な虐待的・残酷な振る舞い」を禁じる条項が注目されたが、通常の不満表明、創作、研究・テストは対象外とする。主な執行手段は会話の終了であり、モデルに人間と同じ権利を認めたという解釈は公式説明を超える。プロパガンダや監視などの制限拡大も報じられた。
\dailyxsource{Hayden Field (@haydenfield)}{https://x.com/haydenfield/status/2108241513159528705}

\section{Anthropic Cyber MissionとOSS Scanner}
Anthropicは重要インフラとオープンソースの防御を支援するCyber Missionを発表した。参加を希望するOSSを定期的に調べる無償のOSS Scannerと、電力・水道などの事業者を対象とするCritical Infrastructure Defense Programが柱となる。スキャナーの報告はモデル生成で人手のレビューを経ないため、重大度の誤判定などがあり得ると公式に注意を促している。防御効果の独立した定量評価は未確認。
\dailyxsource{Anthropic (@AnthropicAI)}{https://x.com/AnthropicAI/status/2108302539498414208}

\section{Anthropic、Genesis Missionに1.5億ドル}
Anthropicは米国のGenesis Missionに今後3年間で1億5,000万ドルを投じ、15以上の連邦機関でClaudeと技術支援を利用できるようにすると発表した。NASA、NIH、NSFなどを含み、研究プロジェクトへの利用クレジット、核融合・量子分野での協力、研修や導入支援を挙げている。科学的発見の加速は同社の期待であり、資金の現金・現物内訳や実際の成果はまだ確認されていない。
\dailyxsource{Anthropic (@AnthropicAI)}{https://x.com/AnthropicAI/status/2108226292235809081}

\section{GPT-6.1 Sol Ultrafastを提供開始}
OpenAI DevelopersはGPT-6.1 Sol UltrafastをAPI、Codex、ChatGPT Workへ順次展開すると投稿した。Sol Standardに対して最大8倍速く、Astraに近い知能を持つという性能表現は公式投稿の主張で、測定条件や独立評価は未確認。ProやEnterpriseなどの対象範囲、料金、地域別の推論レジデンシーについては二次投稿が詳細を述べるが、今回のGrok収集では公式価格表を確認していない。
\dailyxsource{OpenAI Developers (@OpenAIDevs)}{https://x.com/OpenAIDevs/status/2108262812489531498}

\section{StepFun、Step 5 Previewを試験提供}
StepFunはStep 5 PreviewをOpenRouterで提供し、OpenCode、Cline、Nous Research、Kilo Codeなどから1週間無料で試せると告知した。エージェント向けの高性能と低いタスク単価を訴求する一方、パラメータ数、ライセンス、正式なモデルIDや価格は今回の公式投稿から確認できていない。600B規模や10月15日のオープンウェイト公開といった見通しは第三者投稿に限られ、確定情報として扱わない。
\dailyxsource{StepFun (@StepFun\_ai)}{https://x.com/StepFun_ai/status/2108182763002278158}

\section{Claude Scienceによる紫外全天マップ}
Anthropicは天体物理学者Brice MénardがClaude Scienceを用い、既存データセットの統合と統計的推論によって紫外域の全天マップを作成したと紹介した。人手なら数週間かかる作業を数日で進めたという時間短縮は同社の説明である。関連する科学ブログのURL、使用データ、未観測領域を推定した割合、査読状況はGrok収集では確認されていない。
\dailyxsource{Anthropic (@AnthropicAI)}{https://x.com/AnthropicAI/status/2108290395599667700}

\section{エージェントの時間認識を評価}
Maksym Andriushchenkoは、エージェントが指示された作業時間を守れるか、所要時間を予測・回顧できるかを測る研究を紹介した。Astraは時間制約の遵守で優れ、Fable 5.1は4時間指定の作業を1時間未満で終える場合が多いという比較を示したが、いずれも著者の主張である。論文本文や評価条件、図中の数値は今回未取得で、モデル間の一般的な優劣は断定できない。
\dailyxsource{Maksym Andriushchenko (@maksym\_andr)}{https://x.com/maksym_andr/status/2108219551984844851}

\section{FIGS、共感と迎合を分離して評価}
Sahar Abdelnabiは、正確性を優先すると迎合が減る反面、冷たい応答になり得るという問題に対し、FIGSを紹介した。事実への忠実さと相手の感情・状況への適切な配慮を別々に評価し、複数ターンで新証拠による正当な訂正と、圧力に屈した不当な同意を区別する。既存評価の限界を補う提案だが、論文URLや第三者の再現結果は未確認である。
\dailyxsource{Sahar Abdelnabi (@sahar\_abdelnabi)}{https://x.com/sahar_abdelnabi/status/2108270611533312350}

\section{Higgsfield Katana、Claudeで動画編集}
HiggsfieldはClaudeから利用できる動画編集ツールKatana（Claude Motion）を紹介した。参照素材をアップロードし、編集可能なモーショングラフィックスや製品発表向け動画を作る用途を示し、Higgsfield MCP経由で利用できると説明している。具体的な編集範囲や再現性、利用条件については会社側のデモ・説明に基づくもので、独立した品質評価は確認されていない。
\dailyxsource{Higgsfield AI (@higgsfield)}{https://x.com/higgsfield/status/2108294775644585998}

\section{Windows MLがllama.cppに対応}
Windows DeveloperはWindows MLでllama.cppをサポートし、GPU、NPU、CPUを利用してオープンモデルを試せるようにすると投稿した。ローカル推論をWindowsの共通基盤に近づける動きとして注目される。ただしリンク先ブログの日付は10月7日を含み、記事自体は観測窓の直前に公開された可能性がある。本号では10月8日から9日の窓内にあった公式X告知として扱う。
\dailyxsource{Windows Developer (@windowsdev)}{https://x.com/windowsdev/status/2108299049338011699}

\section{Artificial Analysis、Index v5を予告}
Artificial AnalysisはIntelligence Index v5を10月下旬に公開すると予告した。Terminal-Bench Scienceと、非公開データセットを用いる新しいコーディングベンチマークを追加するという。既存のモデル比較指標にエージェント的・実務的な評価を取り込む動きだが、v5の最終的な構成や各モデルの順位は発表前であり、今回の投稿からは判断できない。
\dailyxsource{Artificial Analysis (@ArtificialAnlys)}{https://x.com/ArtificialAnlys/status/2108258926680813846}

\section{米国の科学・技術勲章を授与}
Rapid Response 47は、Sergey Brin、Jensen Huang、Elon Musk、Lisa SuがNational Medal of Scienceを、Michael DellとSatya NadellaがNational Medal of Technology and Innovationを授与されたと動画付きで投稿した。NVIDIA公式もHuangの受賞を祝福し、GPU計算による科学研究への貢献を挙げた。本号では授与に関する窓内の投稿を記録する。
\dailyxsource{Rapid Response 47 (@RapidResponse47)}{https://x.com/RapidResponse47/status/2108258677773967451}

\section{Google、ブラインドスイープ超音波研究}
Googleは母体ヘルスケアへの応用を想定した、ブラインドスイープ超音波とAIに関する研究ブログを紹介した。専門的な撮像操作を簡素化する方向性が注目されるが、公式X投稿には臨床性能の具体的な数値は示されていない。Grok収集ではブログ本文も未読のため、診断精度や実運用の適用範囲はここでは評価しない。
\dailyxsource{Google (@Google)}{https://x.com/Google/status/2108222019980403023}

\section{Grokipedia v0.3、1週間の更新状況}
Grokipediaアカウントは、v0.3公開から1週間で4,400件以上の記事を追加し、Grokによる編集が1日2,200件以上に達したと投稿した。生成AIを使った知識ベースの継続更新を示す数字だが、集計方法、編集の品質、誤り訂正の実績は未確認である。件数は運営側の自己申告として記録する。
\dailyxsource{Grokipedia (@Grokipedia)}{https://x.com/Grokipedia/status/2108313780681875598}

\section{規則遵守に影響するフレーミング}
Mika Okamotoは、LLMやAIアシスタントの規則遵守が説明の仕方、金銭的誘因、社会的圧力によって変化するという研究を紹介した。罰金に触れるだけで遵守や安全性が下がる場合もあるというのは著者の主張で、COLM 2026のAgent Behaviorワークショップなどで発表予定とする。実験設定や効果量は今回未確認であり、全モデルに当てはまるとは言えない。
\dailyxsource{Mika Okamoto (@mikahokamoto)}{https://x.com/mikahokamoto/status/2108241613675966817}

\section{長文脈トークン退避と研究動向}
Stephen Gouldは、将来の文脈を使って冗長なトークンを見極め、退避する手法「Back from the Future」を論文とコード付きで紹介した。同じ観測窓では、Xin Eric Wangがエージェントの好奇心を扱うposition paperのNeurIPS 2026 Spotlight選出を報告し、Charles FryeがLLM Engine Advisorの更新を案内した。後二者は関連研究・ツールの動向として記録し、論文内容の追加検証は行っていない。
\dailyxsource{Stephen Gould (@sgould\_au)}{https://x.com/sgould_au/status/2108316305434804535}

